\section{Source Section II.A: Equilibrium van der Waals Baseline}
\label{sec:vdw-baseline}

\sourceeqrange{Eqs. (2.1)--(2.11)}

\subsection{Purpose and Scope}

This section gives the equilibrium thermodynamic baseline used later in the
companion.  The source begins with a one-component van der Waals fluid and then
adds a square-gradient free-energy term for an equilibrium interface.  The role
of this companion section is to make the local variables, partial derivatives,
and one-dimensional interface bookkeeping explicit before the nonequilibrium
entropy and energy functionals are introduced in Section~\ref{sec:gradient-entropy-energy}.

The scope is equilibrium and local.  We derive the bulk pressure, entropy,
internal-energy, spinodal, and sound-speed identities from the stated Helmholtz
free-energy density.  We then record how the equilibrium gradient term leads to
the chemical-potential and surface-tension formulas used as the capillarity
baseline.  We do not derive the dynamic stress tensor here, do not choose the
Appendix-B printed/scaled branch, and do not make any numerical or PDE
well-posedness claim.

\subsection{Typed Objects and Assumptions}

Let \(n>0\) be the number density and \(T>0\) the absolute temperature.  The
baseline is a local equilibrium model, so the independent variables in this
subsection are \((n,T)\), except in the final one-dimensional interface branch
where \(n=n(x)\) at fixed \(T\).

\begin{center}
\begin{tabular}{@{}lll@{}}
\toprule
Symbol & Type & Role in Section II.A \\
\midrule
\(n\) & scalar density & number density \\
\(T\) & scalar & temperature \\
\(v_0\) & volume parameter & molecular volume \\
\(\epsilon\) & energy parameter & attractive-interaction scale \\
\(k_B\) & constant & Boltzmann constant \\
\(m\) & mass parameter & particle mass \\
\(d\) & integer parameter & space dimension \\
\(\lambda_{\rm th}(T)\) & length & thermal de Broglie wavelength \\
\(f(n,T)\) & energy density & Helmholtz free-energy density \\
\(e(n,T)\) & energy density & internal-energy density \\
\(s(n,T)\) & entropy per particle & source entropy variable \\
\(p(n,T)\) & pressure & bulk thermodynamic pressure \\
\(\mu(n,T)\) & energy per particle & chemical potential \((\partial f/\partial n)_T\) \\
\(T_s(n)\) & temperature & spinodal temperature function \\
\(c\) & speed & adiabatic sound speed \\
\(\gamma_s\) & dimensionless & specific-heat ratio \(C_p/C_V\) \\
\(M(n)\) & coefficient & equilibrium square-gradient free-energy coefficient \\
\bottomrule
\end{tabular}
\end{center}

The local thermodynamic derivatives in this section are classical partial
derivatives.  When a derivative is taken at fixed entropy per particle, the
constraint is written explicitly.  The one-dimensional interface calculation
assumes enough smoothness for \(n_x\) and \(n_{xx}\) to exist and decay or match
the coexisting bulk states at the ends of the interface.

\subsection{Bulk Free Energy and Derived State Functions}

The Helmholtz density in \eqsrc{2.1} has the form
\begin{equation}
  f(n,T)=k_BTn\left[\log\left(\lambda_{\rm th}^d n\right)-1-
  \log(1-v_0n)\right]-\epsilon v_0 n^2.\sourceeqbadge{2.1}
  \label{eq:iia-free-energy}
\end{equation}
The companion treats \(\lambda_{\rm th}^d\) as a temperature-dependent factor
proportional to \(T^{-d/2}\).  At fixed \(T\) it is constant with respect to
\(n\).  The two derivative operations used below are deliberately kept
separate:
\begin{center}
\begin{tabular}{@{}lll@{}}
\toprule
Operation & Held fixed & Differentiation rule \\
\midrule
\(\partial_n\) & \(T\) & \(\partial_n\log(\lambda_{\rm th}^d n)=1/n\) \\
\(\partial_n\) & \(T\) & \(\partial_n[-\log(1-v_0n)]=v_0/(1-v_0n)\) \\
\(\partial_T\) & \(n\) & \(\partial_T\log(\lambda_{\rm th}^d)=-d/(2T)\) \\
\bottomrule
\end{tabular}
\end{center}
No spatial derivative is being taken in this bulk subsection.  The symbol
\(n\) is an independent thermodynamic coordinate, not yet a field \(n(x)\).
For the derivative with respect to \(n\), set
\[
  L(n,T)=\log\left(\frac{\lambda_{\rm th}^d n}{1-v_0n}\right).
\]
Then
\begin{align}
  \left(\frac{\partial}{\partial n}n[L(n,T)-1]\right)_T
  &=L-1+n\left(\frac{1}{n}+\frac{v_0}{1-v_0n}\right) \\
  &=L+\frac{v_0n}{1-v_0n}.
  \label{eq:iia-chemical-potential-log-step}
\end{align}
This is a companion derivation step, not an additional source equation.  It
makes the fixed-temperature derivative in the chemical potential explicit:
\begin{equation}
  \mu=\left(\frac{\partial f}{\partial n}\right)_T
  =k_BT\left[\log\left(\frac{\lambda_{\rm th}^d n}{1-v_0n}\right)
  +\frac{v_0n}{1-v_0n}\right]-2\epsilon v_0n .
  \label{eq:iia-chemical-potential}
\end{equation}
The pressure follows from the Euler relation for the Helmholtz density,
\begin{equation}
  p=n\mu-f.
  \label{eq:iia-pressure-euler}
\end{equation}
The step-by-step calculation is a cancellation ledger.  First multiply
\(\mu\) by \(n\), then subtract the entire Helmholtz density; only after this
subtraction should the attractive term be simplified.  Written in three groups,
\[
  n\mu-f
  =\underbrace{k_BTnL-k_BTn(L-1)}_{\text{ideal logarithmic part}}
  +\underbrace{k_BTn\frac{v_0n}{1-v_0n}}_{\text{excluded-volume derivative}}
  +\underbrace{(-2\epsilon v_0n^2+\epsilon v_0n^2)}_{\text{attraction}} .
\]
The logarithmic terms cancel when this relation is expanded:
\begin{align}
  n\mu-f
  &=k_BTn\left[L+\frac{v_0n}{1-v_0n}-(L-1)\right]
    -\epsilon v_0n^2 \\
  &=k_BTn\left(1+\frac{v_0n}{1-v_0n}\right)-\epsilon v_0n^2 \\
  &=\frac{nk_BT}{1-v_0n}-\epsilon v_0n^2.
  \label{eq:iia-pressure-cancellation-step}
\end{align}
Thus \eqref{eq:iia-free-energy} and \eqref{eq:iia-chemical-potential}
give exactly the pressure formula in \eqsrc{2.4},
\begin{equation}
  p(n,T)=\frac{nk_BT}{1-v_0n}-\epsilon v_0n^2 .\sourceeqbadge{2.4}
  \label{eq:iia-pressure}
\end{equation}
This is a thermodynamic identity for the local bulk model, not a mechanical
stress law.

For later checking it is useful to isolate the two cancellations in the bulk
state-function chain.  First,
\begin{align}
  L+\frac{v_0n}{1-v_0n}-(L-1)
  &=1+\frac{v_0n}{1-v_0n}
    =\frac{1}{1-v_0n}, \\
  -2\epsilon v_0n^2-(-\epsilon v_0n^2)
  &=-\epsilon v_0n^2 .
  \label{eq:iia-pressure-linebyline-ledger}
\end{align}
Second, every logarithmic term in \(e=f+Tns\) cancels before the attraction
term is carried down unchanged.  These are companion audit lines: they explain
why the source pressure and energy formulas follow from the Helmholtz density,
but they are not additional source equations.

The entropy per particle is recovered by differentiating \(f\) with respect to
\(T\) at fixed \(n\):
\begin{equation}
  ns=-\left(\frac{\partial f}{\partial T}\right)_n .
  \label{eq:iia-entropy-definition}
\end{equation}
Because \(\lambda_{\rm th}^d\propto T^{-d/2}\),
\begin{equation}
  \left(\frac{\partial L}{\partial T}\right)_n=-\frac{d}{2T}.
  \label{eq:iia-thermal-wavelength-derivative}
\end{equation}
Therefore the fixed-density derivative is
\begin{align}
  \left(\frac{\partial f}{\partial T}\right)_n
  &=k_Bn(L-1)+k_BTn\left(-\frac{d}{2T}\right) \\
  &=k_Bn\left(L-1-\frac{d}{2}\right),
  \label{eq:iia-entropy-temperature-derivative-step}
\end{align}
and \eqref{eq:iia-entropy-definition} gives the entropy formula \eqsrc{2.3},
\begin{equation}
  s(n,T)=-k_B\log\left(\frac{\lambda_{\rm th}^d n}{1-v_0n}\right)
  +\frac{k_B(d+2)}{2} .\sourceeqbadge{2.3}
  \label{eq:iia-entropy}
\end{equation}
The internal-energy density follows from \(e=f+Tns\).  Substituting the
just-derived entropy makes the cancellation visible:
\begin{align}
  f+Tns
  &=k_BTn(L-1)-\epsilon v_0n^2
    +k_BTn\left(-L+1+\frac{d}{2}\right) \\
  &=\frac{d}{2}nk_BT-\epsilon v_0n^2.
  \label{eq:iia-internal-energy-cancellation-step}
\end{align}
This yields \eqsrc{2.2},
\begin{equation}
  e(n,T)=\frac{d}{2}nk_BT-\epsilon v_0n^2 .\sourceeqbadge{2.2}
  \label{eq:iia-internal-energy}
\end{equation}
The check fixes the meaning of \(s\): it is an entropy per particle, so the
entropy density is \(ns\).

\subsection{Coexistence Constants and Profile Terminology}

At fixed temperature below the critical point, the equilibrium interface branch
uses two bulk densities, denoted here by \(n_\alpha\) and \(n_\beta\).  They
are not chosen by the gradient calculation alone.  They are the bulk endpoints
of the coexistence construction: the chemical potential and pressure take the
same constant values in both bulk phases,
\begin{equation}
  \mu(n_\alpha,T)=\mu(n_\beta,T)=\mu_{\rm cx}(T),\qquad
  p(n_\alpha,T)=p(n_\beta,T)=p_{\rm cx}(T).
  \label{eq:iia-coexistence-bulk-conditions}
\end{equation}
Equation \eqref{eq:iia-coexistence-bulk-conditions} is the companion's
notation for the coexistence branch used by source Eqs.~(2.9)--(2.11).  It is a
bulk thermodynamic condition, not a dynamic boundary condition and not a
numerical prescription.

The one-dimensional object used below is an equilibrium density profile,
\(n_*(x)\), connecting the two bulk limits through a diffuse interface,
\begin{equation}
  \lim_{x\to-\infty}n_*(x)=n_\alpha,\qquad
  \lim_{x\to+\infty}n_*(x)=n_\beta,\qquad
  \lim_{|x|\to\infty}\partial_x n_*(x)=0 .
  \label{eq:iia-equilibrium-profile-limits}
\end{equation}
The word ``profile'' is used in this restricted sense: it is a spatial
stationary field at fixed \(T\), not a time trajectory and not a sharp-interface
jump condition.  These endpoint statements explain how the integration
constants in the pressure and surface-tension formulas are fixed.


\begin{figure}[htbp]
\centering
\resizebox{0.82\linewidth}{!}{\begin{tikzpicture}[
  font=\scriptsize,
  axis/.style={-{Latex[length=2mm]}, line width=0.45pt},
  curve/.style={line width=0.9pt},
  guide/.style={gray!65, dashed, line width=0.35pt},
  box/.style={draw=black, line width=0.4pt, rounded corners=2pt,
    align=center, text width=24mm, fill=gray!8, minimum height=7mm},
  note/.style={align=center, text width=18mm}
]
\fill[gray!10] (-0.1,0.15) rectangle (1.75,2.95);
\fill[gray!18] (4.25,0.15) rectangle (6.1,2.95);
\draw[axis] (-0.2,0.2) -- (6.35,0.2) node[right] {$x$};
\draw[axis] (0,0.0) -- (0,3.18) node[above] {$n_*(x)$};
\draw[guide] (0,0.85) -- (6.0,0.85);
\draw[guide] (0,2.45) -- (6.0,2.45);
\node[left] at (0,0.85) {$n_\alpha$};
\node[left] at (0,2.45) {$n_\beta$};
\node[note] at (0.86,0.43) {bulk $\alpha$};
\node[note] at (5.18,2.72) {bulk $\beta$};
\draw[curve] plot[smooth] coordinates {
  (0.2,0.86) (1.2,0.88) (2.1,1.08) (3.0,1.65) (3.9,2.25) (4.8,2.43) (5.8,2.45)
};
\node[box] at (2.75,3.06) {diffuse profile\\not a sharp jump};
\draw[-{Latex[length=2mm]}, line width=0.4pt] (2.75,2.81) -- (3.02,1.72);
\node[box] at (3.05,-0.42) {coexistence fixes\\$\mu_\mathrm{cx}$ and $p_\mathrm{cx}$};
\draw[-{Latex[length=2mm]}, line width=0.4pt] (2.25,-0.2) -- (1.05,0.78);
\draw[-{Latex[length=2mm]}, line width=0.4pt] (3.85,-0.2) -- (4.9,2.36);
\end{tikzpicture}
}
\caption{Navigation diagram for the one-dimensional equilibrium density
profile used in the Section II.A interface calculation.  The curve represents a
diffuse stationary profile joining two bulk coexistence limits; it is not a
sharp-interface jump condition and not a time trajectory.}
\label{fig:coexistence-profile-navigation}
\end{figure}


\subsection{Spinodal Temperature, Sound Speed, and Heat-Capacity Ratio}

Differentiate \eqref{eq:iia-pressure} at fixed \(T\):
\begin{equation}
  \left(\frac{\partial p}{\partial n}\right)_T
  =\frac{k_BT}{(1-v_0n)^2}-2\epsilon v_0 n .
  \label{eq:iia-isothermal-pressure-derivative}
\end{equation}
The spinodal temperature recorded in \eqsrc{2.7} is the value of \(T\) for
which this isothermal derivative vanishes,
\begin{equation}
  T_s(n)=\frac{2\epsilon v_0n(1-v_0n)^2}{k_B} .\sourceeqbadge{2.7}
  \label{eq:iia-spinodal-temperature}
\end{equation}
Thus
\begin{equation}
  \left(\frac{\partial p}{\partial n}\right)_T
  =\frac{k_B(T-T_s)}{(1-v_0n)^2} .
  \label{eq:iia-isothermal-derivative-spinodal-form}
\end{equation}
To locate the maximum, write \(u=v_0n\).  Apart from the positive constant
\(2\epsilon/k_B\), the spinodal function is \(u(1-u)^2\), and
\[
  \frac{\dd}{\dd u}u(1-u)^2=(1-u)(1-3u).
\]
The nontrivial maximum in the physical interval therefore occurs at
\(u=1/3\), or \(n=1/(3v_0)\), giving the familiar critical values recorded
after \eqsrc{2.7},
\begin{equation}
  n_c=\frac{1}{3v_0},\qquad T_c=\frac{8\epsilon}{27k_B}.
  \label{eq:iia-critical-values}
\end{equation}

For the adiabatic derivative, impose \(\dd s=0\).  The calculation uses
the chain rule along the curve \(T=T(n)\) on which the entropy per particle is
constant:
\[
  \left(\frac{\dd p}{\dd n}\right)_{s}
  =\left(\frac{\partial p}{\partial n}\right)_{T}
  +\left(\frac{\partial p}{\partial T}\right)_{n}
   \left(\frac{\partial T}{\partial n}\right)_{s}.
\]
The three factors on the right have different types: the first two are local
thermodynamic partial derivatives, while the last is the slope of the
constant-entropy curve in the \((n,T)\) plane.  From
\eqref{eq:iia-entropy},
\begin{equation}
  \left(\frac{\partial s}{\partial n}\right)_T
  =-\frac{k_B}{n(1-v_0n)},\qquad
  \left(\frac{\partial s}{\partial T}\right)_n=\frac{dk_B}{2T}.
  \label{eq:iia-entropy-partials}
\end{equation}
The fixed-entropy differential is
\[
  0=\dd s=\left(\frac{\partial s}{\partial n}\right)_T\dd n
  +\left(\frac{\partial s}{\partial T}\right)_n\dd T,
\]
so the derivative is obtained by solving the scalar linear equation
\begin{equation}
  \left(\frac{\partial s}{\partial T}\right)_n
  \left(\frac{\partial T}{\partial n}\right)_s
  =-\left(\frac{\partial s}{\partial n}\right)_T .
  \label{eq:iia-fixed-entropy-linear-solve}
\end{equation}
Thus
\begin{equation}
  \left(\frac{\partial T}{\partial n}\right)_s
  =-\frac{(\partial s/\partial n)_T}{(\partial s/\partial T)_n}
  =\frac{2T}{dn(1-v_0n)} .
  \label{eq:iia-adiabatic-temperature-derivative}
\end{equation}
Using \((\partial p/\partial T)_n=nk_B/(1-v_0n)\), the adiabatic pressure
derivative is
\begin{align}
  \left(\frac{\partial p}{\partial n}\right)_s
  &=\left(\frac{\partial p}{\partial n}\right)_T
    +\left(\frac{\partial p}{\partial T}\right)_n
     \left(\frac{\partial T}{\partial n}\right)_s \\
  &=\frac{k_BT}{(1-v_0n)^2}
    \left(1+\frac{2}{d}-\frac{T_s}{T}\right).
  \label{eq:iia-adiabatic-pressure-derivative}
\end{align}
Dividing by the particle mass \(m\) gives the derivative-based sound-speed
formula
\begin{equation}
  c_{\rm der}=\frac{1}{1-v_0n}\left(\frac{k_BT}{m}\right)^{1/2}
    \left(1+\frac{2}{d}-\frac{T_s}{T}\right)^{1/2} .
    \companionbranchbadge{derived 2.5 branch}
  \label{eq:iia-sound-speed}
\end{equation}
The badge identifies this as the companion's derivative-consistent branch,
not as a transcription of the printed formula.  The page image for printed
source Eq.~(2.5) displays the same square-root factor
but without the prefactor \((1-v_0n)^{-1}\).  The denominator is forced by
\eqref{eq:iia-adiabatic-pressure-derivative} and is also present in the
Section III scaled coefficient, source Eq.~(3.9).  This companion therefore
keeps the derivative-consistent expression in \eqref{eq:iia-sound-speed} and
records the printed source Eq.~(2.5) line as a source-fidelity discrepancy, not
as a formal erratum and not as a change to any simulation branch.
The heat-capacity ratio is the ratio of the adiabatic and isothermal pressure
responses,
\begin{equation}
  \gamma_s=\frac{C_p}{C_V}
  =\frac{(\partial p/\partial n)_s}{(\partial p/\partial n)_T}
  =1+\frac{2}{d(1-T_s/T)} .\sourceeqbadge{2.6}
  \label{eq:iia-heat-capacity-ratio}
\end{equation}
Equation \eqref{eq:iia-heat-capacity-ratio} writes \eqsrc{2.6} in an
unambiguous fraction form.

\subsection{Equilibrium Gradient Free Energy}

The equilibrium capillarity branch begins with the square-gradient density in
\eqsrc{2.8},
\begin{equation}
  f_{\rm gra}=\frac{1}{2}M(n)|\grad n|^2 .\sourceeqbadge{2.8}
  \label{eq:iia-gradient-free-energy}
\end{equation}
For a planar interface, take \(n=n(x)\) and fixed \(T\).  Vary the
one-dimensional functional
\begin{equation}
  F_{1d}[n]=\int \left[f(n,T)+\frac{1}{2}M(n)n_x^2\right]\dd x
  \label{eq:iia-one-dimensional-functional}
\end{equation}
by \(n_\varepsilon=n+\varepsilon\eta\), where \(\eta\) is a smooth test
variation.  The first-order term is
\begin{align}
  \left.\frac{\dd}{\dd\varepsilon}F_{1d}[n_\varepsilon]\right|_{\varepsilon=0}
  &=\int \left[\mu\eta+\frac{1}{2}M_n n_x^2\eta
    +M n_x\eta_x\right]\dd x \\
  &=\int \left[\mu-Mn_{xx}-\frac{1}{2}M_n n_x^2\right]\eta\,\dd x
    +\left[M n_x\eta\right]_{\partial}.
  \label{eq:iia-gradient-variation-step}
\end{align}
The integration-by-parts sign in the second line comes from the one-dimensional
product rule
\begin{equation}
  \frac{\dd}{\dd x}(M n_x\eta)
  =M n_x\eta_x+(M n_{xx}+M_n n_x^2)\eta .
  \label{eq:iia-gradient-product-rule-linebyline}
\end{equation}
Hence
\begin{equation}
  \int M n_x\eta_x\,\dd x
  =[M n_x\eta]_{\partial}
  -\int (M n_{xx}+M_n n_x^2)\eta\,\dd x .
  \label{eq:iia-gradient-ibp-linebyline}
\end{equation}
Here \(M_n=(\partial M/\partial n)_T\) is a thermodynamic partial derivative,
not a spatial derivative.  The boundary term should be read before any
endpoint assumption is imposed:
\begin{center}
\begin{tabular}{@{}lll@{}}
\toprule
Branch & Assumption on \(\eta\) & Consequence for \([Mn_x\eta]_{\partial}\) \\
\midrule
fixed endpoints & \(\eta=0\) at both ends & term vanishes \\
free endpoint & \(\eta\) arbitrary & natural boundary data are required \\
coexisting bulk limit & \(n_x\to0\) at infinity & term vanishes if the decay is strong enough \\
\bottomrule
\end{tabular}
\end{center}
With endpoint variations fixed, the boundary term
vanishes and the Euler--Lagrange derivative is
\begin{equation}
  \frac{\delta F_{1d}}{\delta n}
  =\mu(n,T)-M n_{xx}-\frac{1}{2}M_n n_x^2,
  \qquad M_n=\left(\frac{\partial M}{\partial n}\right)_T .
  \label{eq:iia-gradient-euler-lagrange}
\end{equation}
The coexistence chemical potential \(\mu_{\rm cx}(T)\) is the constant value in
\eqref{eq:iia-coexistence-bulk-conditions}.  A stationary profile has the same
Euler--Lagrange derivative at every \(x\), so rearranging gives the
chemical-potential interface relation in \eqsrc{2.9}, in component notation,
\begin{equation}
  \mu-\mu_{\rm cx}=M n_{xx}+\frac{1}{2}M_n n_x^2
  =\frac{M}{2}n_{xx}+\frac{\dd}{\dd x}\left(\frac{M}{2}n_x\right).\sourceeqbadge{2.9}
  \label{eq:iia-interface-chemical-potential}
\end{equation}
The second equality is often the safest way to read the source notation when
\(M=M(n)\).

\subsection{Fixed-Temperature Route to the Interface Pressure Relation}

The source next uses the fixed-temperature thermodynamic identity
\(\dd p=n\,\dd\mu\) to obtain \eqsrc{2.10}.  This is a local
one-dimensional calculation along the equilibrium profile, not a dynamic stress
balance.  Let
\begin{equation}
  A=M n_{xx}+\frac{1}{2}M_x n_x
   =\frac{M}{2}n_{xx}+\frac{\dd}{\dd x}\left(\frac{M}{2}n_x\right),
  \label{eq:iia-interface-force-abbrev}
\end{equation}
so that \eqref{eq:iia-interface-chemical-potential} says
\(\mu-\mu_{\rm cx}=A\).  Because the temperature is fixed and
\(\mu_{\rm cx}\) is constant along the profile,
\begin{equation}
  \frac{\dd}{\dd x}(p-p_{\rm cx})
  =n\frac{\dd}{\dd x}(\mu-\mu_{\rm cx})
  =nA_x .
  \label{eq:iia-dp-ndmu-route}
\end{equation}
Now define the candidate pressure excess
\begin{equation}
  P=nA-\frac{1}{2}M n_x^2 .
  \label{eq:iia-pressure-candidate}
\end{equation}
Its derivative is
\begin{align}
  P_x
  &=n_xA+nA_x-\frac{1}{2}M_xn_x^2-Mn_xn_{xx} \\
  &=nA_x+
    \underbrace{n_x\left(Mn_{xx}+\frac{1}{2}M_xn_x\right)
    -\frac{1}{2}M_xn_x^2-Mn_xn_{xx}}_{0} \\
  &=nA_x .
  \label{eq:iia-pressure-candidate-derivative}
\end{align}
Thus \(P\) and \(p-p_{\rm cx}\) have the same derivative.  Their integration
constant is fixed by the coexistence endpoint conditions
\eqref{eq:iia-coexistence-bulk-conditions}--\eqref{eq:iia-equilibrium-profile-limits}:
in the two bulk limits \(n_x=n_{xx}=0\) and \(p=p_{\rm cx}\), so
\(P=p-p_{\rm cx}\).  Rewriting \(P\) in the compact source form gives
\begin{equation}
  p-p_{\rm cx}
  =\frac{M}{2}n n_{xx}
   +\frac{\dd}{\dd x}\left(\frac{M}{2}n n_x\right)
   -M n_x^2 .\sourceeqbadge{2.10}
  \label{eq:iia-interface-pressure-relation}
\end{equation}
Expanding the total derivative in \eqref{eq:iia-interface-pressure-relation}
recovers the equivalent form \(Mn n_{xx}+\frac{1}{2}nM_xn_x-\frac{1}{2}Mn_x^2\).
This expression is the equilibrium pressure bridge to the Section II.D pressure
tensor; it is not a decision about the Appendix-B printed/scaled branch.

The surface-tension formula in \eqsrc{2.11} is the positive interfacial contribution.
Its provenance is the same square-gradient equilibrium branch, not a separate
dynamic law.  Let
\(\omega(n,T)=f(n,T)-\mu_{\rm cx}(T)n\) be the bulk grand-potential density at
the coexistence chemical potential.  Along the stationary profile,
\begin{equation}
  \frac{\dd}{\dd x}\bigl(\omega-\omega_{\rm cx}\bigr)
  =\bigl(\mu-\mu_{\rm cx}\bigr)n_x
  =\frac{\dd}{\dd x}\left(\frac{1}{2}M n_x^2\right),
  \label{eq:iia-grand-potential-first-integral}
\end{equation}
where the last equality uses \eqref{eq:iia-interface-chemical-potential} times
\(n_x\).  The endpoint limits in \eqref{eq:iia-equilibrium-profile-limits}
make both sides vanish in the bulk phases, so
\begin{equation}
  \omega-\omega_{\rm cx}=\frac{1}{2}M n_x^2 .
  \label{eq:iia-grand-potential-excess}
\end{equation}
The interfacial excess of bulk grand potential plus square-gradient energy is
therefore
\begin{equation}
  \gamma(T)=\int_{-\infty}^{\infty}\left[
    \omega-\omega_{\rm cx}+\frac{1}{2}M(n(x))n_x^2
  \right]\dd x
  =\int_{-\infty}^{\infty}M(n(x))n_x^2\dd x .\sourceeqbadge{2.11}
  \label{eq:iia-surface-tension}
\end{equation}
For a monotone profile oriented so that \(n_x>0\), use
\(\dd n=n_x\dd x\).  The first integral
\eqref{eq:iia-grand-potential-excess} gives
\(n_x=[2(\omega-\omega_{\rm cx})/M]^{1/2}\), and hence
\begin{align}
  \gamma(T)
  &=\int_{n_\alpha}^{n_\beta}M n_x\,\dd n \nonumber\\
  &=\int_{n_\alpha}^{n_\beta}
    \left[2M(n,T)\bigl(\omega(n,T)-\omega_{\rm cx}(T)\bigr)\right]^{1/2}
    \dd n .
  \label{eq:iia-surface-tension-density-change}
\end{align}
If the profile orientation is reversed, the integration limits reverse with
it; the physical surface tension remains positive.  This change of variables
therefore belongs to the monotone equilibrium branch and is not available for
an arbitrary nonmonotone field.
This derivation records the square-gradient first-integral route to the source
formula.  The source notes that Eqs.~(2.9)--(2.11) still apply when \(M\)
depends on \(n\).  In this companion, all such statements remain equilibrium
one-dimensional statements until Section II.D introduces the dynamic pressure
tensor.

\subsection{Connection to Later Sections}

Section~\ref{sec:gradient-entropy-energy} replaces the equilibrium coefficient
\(M\) by the nonequilibrium combination \(M=CT+K\).  Appendix~\ref{app:reversible-stress-derivation}
and Section~\ref{sec:hydrodynamic-equations} then use the same capillary
coefficient inside the reversible stress tensor.  Appendix~\ref{app:scaled-equations-plan}
uses the simulation branch \(K=0\), constant \(C\), and \(M=CT\) when rescaling
Appendix-B equations.

The inspected Korteweg 1901 source provides historical and mechanical-stress
context, but it does not replace Onuki as the authority for the equilibrium
free-energy formulas in this section.  Appendix~
\ref{app:korteweg-primary-crosswalk} records the source distinction and the
restricted tensor map separately.

\subsection{Verification Checklist}

The finite verification layer for this section checks:

\begin{itemize}
\item \(p=n\mu-f\) gives \eqref{eq:iia-pressure};
\item \(s=-(1/n)(\partial f/\partial T)_n\) gives \eqref{eq:iia-entropy};
\item \(e=f+Tns\) gives \eqref{eq:iia-internal-energy};
\item \((\partial p/\partial n)_T\) factors through \(T_s\);
\item the adiabatic derivative gives the derivative-consistent sound-speed
  expression, including the excluded-volume denominator flagged above;
\item \(\gamma_s\) equals the ratio of adiabatic to isothermal pressure
  derivatives;
\item the one-dimensional gradient Euler--Lagrange derivative gives
  \eqref{eq:iia-interface-chemical-potential};
\item the monotone change of variables in
  \eqref{eq:iia-surface-tension-density-change} follows from the same first
  integral and introduces no independent surface-tension law;
\item the fixed-temperature route \(\dd p=n\,\dd\mu\) gives
  \eqref{eq:iia-interface-pressure-relation};
\item the surface-tension integrand derivative is consistent with the
  equilibrium first-integral bookkeeping.
\end{itemize}

These checks are algebraic and differential identities.  They do not prove the
physical equation of state, the second law, the dynamic model, numerical
results, or PDE well-posedness.

The executable mirrors for the displayed Section II.A calculations are:
\begin{center}
\small
\begin{tabular}{@{}l@{}}
\texttt{scripts/python/section\_iia\_vdw\_baseline\_sympy.py}\\
\texttt{scripts/wolfram/section\_iia\_vdw\_baseline\_checks.wls}\\
\texttt{tests/test\_section\_iia\_vdw\_baseline.py}
\end{tabular}
\end{center}
In particular, the fixed-temperature pressure route in
\eqref{eq:iia-dp-ndmu-route}--\eqref{eq:iia-interface-pressure-relation} is
checked as a finite one-dimensional identity, not as a dynamic stress law.

\subsection*{Tagged verification complement}

\OnukiLinkIIA{The tagged Section II.A verification complement} gives the hand
route, executable commands, negative control, and finite scope for this local
thermodynamic calculation.

\subsection{Open Issues Carried Forward}

\begin{itemize}
\item The pressure expression in \eqsrc{2.10} is now derived here as the
  fixed-temperature pressure bridge, but it should still be read together with
  the Section II.D pressure tensor before it is used in Appendix-B branch analysis.
\item Appendix-B printed-B3 versus dimensional-source scaling remains separate
  from this equilibrium baseline.
\item The unpublished simulation-code branch remains unresolved.
\end{itemize}
